% !TeX program = xelatex
\documentclass[11pt,a4paper]{ctexart}
\usepackage[margin=2.5cm]{geometry}
\usepackage{amsmath,amssymb}
\setlength{\parindent}{2em}
\setlength{\parskip}{0.4em}
\newcommand{\supp}{\operatorname{supp}}
\title{数学分析课堂笔记}
\author{}
\date{}
\begin{document}
\maketitle

\section*{lecture007}
\subsection*{回顾：一个简单版本的单位分解定理}

设 $K\subset\mathbb R^n$ 为紧集，$\{U_\alpha\}_{\alpha\in\Lambda}$
是 $K$ 的开覆盖。则存在有限子覆盖 $\{U_i\}_{i=1}^k$，以及
$\{\varphi_i\}_{i=1}^k\subset C^\infty(\mathbb R^n)$，使得
\[
 \supp\varphi_i\subset U_i,\qquad 0\leq\varphi_i\leq1,
 \qquad \sum_{i=1}^k\varphi_i=1\quad\text{在 }K\text{ 上}.
\]

\subsection*{Step 1：缩小覆盖}
取开集 $V_1,\ldots,V_k$，使得
\[
 \overline{V_i}\subset U_i,\qquad K\subset\bigcup_{i=1}^kV_i.
\]

\subsection*{Step 2：构造截断函数}
由 $\overline{V_i}\subset U_i$，取 $\psi_i\in C_0^\infty(\mathbb R^n)$，使得
\[
 0\leq\psi_i\leq1,\qquad \supp\psi_i\subset U_i,
 \qquad \psi_i|_{V_i}=1.
\]

\subsection*{Step 3：归一化}
小问题：不能保证 $\sum_j\psi_j$ 在全局都大于 $0$，所以不能在整个
$\mathbb R^n$ 上直接定义 $\psi_i/\sum_j\psi_j$。

\paragraph{Solution 1.}
令
\[
 \Omega=\left\{x\in\mathbb R^n:\sum_{j=1}^k\psi_j(x)>0\right\}.
\]
这是包含 $K$ 的开集。选取截断函数 $\eta$，使
$\eta|_K=1$、$\supp\eta\subset\Omega$，并定义
\[
 \varphi_i(x)=
 \begin{cases}
 \displaystyle\eta(x)\frac{\psi_i(x)}{\sum_{j=1}^k\psi_j(x)},&x\in\Omega,\\[6pt]
 0,&x\in\mathbb R^n\setminus\Omega.
 \end{cases}
\]
在 $\Omega$ 中局部光滑；在 $\mathbb R^n\setminus\Omega$ 的每一点附近，
$\eta$ 恒为 $0$，因此零延拓也是光滑的。

\paragraph{Solution 2.}
令
\[
 \varphi_1=\psi_1,\qquad
 \varphi_2=(1-\psi_1)\psi_2,\qquad\ldots,\qquad
 \varphi_k=\psi_k\prod_{j=1}^{k-1}(1-\psi_j).
\]
这些函数均光滑，且 $0\leq\varphi_i\leq1$、
$\supp\varphi_i\subset\supp\psi_i\subset U_i$。由逐项相消，
\[
 \sum_{i=1}^k\varphi_i
 =1-\prod_{i=1}^k(1-\psi_i).
\]
对任意 $x\in K$，存在 $\ell$ 使 $x\in V_\ell$，于是 $\psi_\ell(x)=1$，
乘积为 $0$，故 $\sum_{i=1}^k\varphi_i(x)=1$。这就得到所需的单位分解。

\subsection*{考虑含参积分}
考虑
\[
 F(\lambda)=\int_a^b f(x)e^{\lambda S(x)}\,dx,
 \qquad f,S\text{ 连续}.
\]
例如，作代换 $x=\lambda t$，有
\[
 \lim_{\lambda\to\infty}\Gamma(\lambda+1)
 =\lim_{\lambda\to\infty}\int_0^\infty x^\lambda e^{-x}\,dx
 =\lim_{\lambda\to\infty}\lambda^{\lambda+1}
   \int_0^\infty e^{\lambda(\ln t-t)}\,dt.
\]

\subsection*{局部化引理：乘法形式}
设 $f,S\in C([a,b])$，$x_0$ 是 $S$ 的唯一全局最大值点，即
\[
 S(x)<S(x_0)\qquad(x\ne x_0),
\]
且 $f(x_0)\ne0$。则对 $x_0$ 的任意充分小的固定相对开邻域
$U\subset[a,b]$，存在 $c>0$，使得当 $\lambda\to+\infty$ 时，
\[
 \boxed{
 \int_a^b f(x)e^{\lambda S(x)}\,dx
 =\left(\int_U f(x)e^{\lambda S(x)}\,dx\right)
   \left(1+O(e^{-c\lambda})\right).
 }
\]

\paragraph{证明.}
记 $M=S(x_0)$。由连续性和唯一全局最大值的条件，紧集
$[a,b]\setminus U$ 上存在一致间隙 $\delta>0$，使
\[
 S(x)\le M-\delta\qquad(x\in[a,b]\setminus U).
\]
因此外部积分满足
\[
 \left|\int_{[a,b]\setminus U}f(x)e^{\lambda S(x)}\,dx\right|
 \le (b-a)\|f\|_\infty e^{\lambda(M-\delta)}.
\]
取 $U$ 足够小，使 $f$ 在 $U$ 上同号；再取更小的相对邻域
$V\subset U$，使 $|V|>0$，且
\[
 |f(x)|\ge m>0,\qquad S(x)\ge M-\frac\delta2
 \qquad(x\in V).
\]
于是
\[
 \left|\int_U f(x)e^{\lambda S(x)}\,dx\right|
 =\int_U |f(x)|e^{\lambda S(x)}\,dx
 \ge m|V|e^{\lambda(M-\delta/2)}.
\]
将外部积分除以局部积分，得到
\[
 \frac{\displaystyle\int_{[a,b]\setminus U}f(x)e^{\lambda S(x)}\,dx}
 {\displaystyle\int_U f(x)e^{\lambda S(x)}\,dx}
 =O(e^{-\lambda\delta/2}),
\]
故可取 $c=\delta/2$。若 $U=[a,b]$，结论直接成立。

因此积分的主要部分只与最大值点附近有关。这里要求唯一的
\emph{全局}最大值点，仅有严格局部极大值不够；$f(x_0)\ne0$
保证充分小邻域内没有正负相消。此引理不要求极大值点非退化。

\paragraph{RMK（局部化引理的加法形式）.}
设 $f,S\in C([a,b])$，$x_0\in[a,b]$ 是 $S$ 的唯一全局最大值点，
记 $M=S(x_0)$。则对 $x_0$ 的任意固定相对开邻域 $U\subset[a,b]$，
存在 $\delta>0$，使得当 $\lambda\to+\infty$ 时，
\[
 \boxed{
 \int_a^b f(x)e^{\lambda S(x)}\,dx
 =\int_U f(x)e^{\lambda S(x)}\,dx
 +O\!\left(e^{\lambda(M-\delta)}\right).
 }
\]
此形式不要求 $f(x_0)\ne0$，也不要求极大值点非退化。

\paragraph{RMK 2（无穷区间上的加法局部化）.}
设 $g,S$ 是 $[0,+\infty)$ 上的连续函数，$S\le M$，
$U$ 是最大值点附近的固定相对开邻域。假设存在 $\delta>0$，使
\[
 S(x)\le M-\delta\qquad(x\in[0,+\infty)\setminus U),
\]
且存在固定的 $\lambda_0\ge0$，使
\[
 \int_0^\infty |g(x)|e^{\lambda_0S(x)}\,dx<\infty.
\]
则当 $\lambda\to+\infty$ 时，
\[
 \int_0^\infty g(x)e^{\lambda S(x)}\,dx
 =\int_U g(x)e^{\lambda S(x)}\,dx
 +O\!\left(e^{\lambda(M-\delta)}\right).
\]
\textbf{简要证明：}记 $U^c=[0,+\infty)\setminus U$。对
$\lambda\ge\lambda_0$，有
\begin{align*}
 \left|\int_{U^c}g(x)e^{\lambda S(x)}\,dx\right|
 &\le\int_{U^c}|g(x)|e^{\lambda_0S(x)}
                 e^{(\lambda-\lambda_0)S(x)}\,dx\\
 &\le e^{(\lambda-\lambda_0)(M-\delta)}
       \int_{U^c}|g(x)|e^{\lambda_0S(x)}\,dx
 =O\!\left(e^{\lambda(M-\delta)}\right).
\end{align*}
因此不必要求 $g$ 本身可积，只需它在某个固定指数权重下绝对可积。
无穷区间上，唯一全局最大值并不自动保证远处与最大值有一致间隙，
故上述间隙条件需明确给出。

\subsection*{Laplace 方法：非严格的解释}
设 $f,S\in C([a,b])$，$x_0\in(a,b)$ 是 $S$ 的唯一全局最大值点，
$S$ 在 $x_0$ 附近二阶连续可微，且
\[
 f(x_0)\ne0,\qquad S'(x_0)=0,\qquad S''(x_0)<0.
\]
这里 $S''(x_0)<0$ 表示极大值点非退化。记
$A=-S''(x_0)>0$，取固定的充分小 $\varepsilon>0$。
由局部化引理，只需考虑 $x_0$ 附近。形式上，用
\[
 f(x)\approx f(x_0),\qquad
 S(x)\approx S(x_0)+S'(x_0)(x-x_0)
             +\frac{S''(x_0)}2(x-x_0)^2,
\]
得到
\begin{align*}
 F(\lambda)
 &\approx\int_{x_0-\varepsilon}^{x_0+\varepsilon}
 f(x_0)\exp\!\left\{\lambda\left[
 S(x_0)+S'(x_0)(x-x_0)+\frac{S''(x_0)}2(x-x_0)^2
 \right]\right\}\,dx\\
 &=f(x_0)e^{\lambda S(x_0)}
 \int_{x_0-\varepsilon}^{x_0+\varepsilon}
 \exp\!\left\{-\frac{\lambda A}{2}(x-x_0)^2\right\}\,dx.
\end{align*}
令
\[
 y=\sqrt{\frac{\lambda A}{2}}(x-x_0),
 \qquad dx=\sqrt{\frac{2}{\lambda A}}\,dy,
\]
则上面的二次近似积分恰好等于
\[
 f(x_0)e^{\lambda S(x_0)}\sqrt{\frac{2}{\lambda A}}
 \int_{-\varepsilon\sqrt{\lambda A/2}}^{\varepsilon\sqrt{\lambda A/2}}
 e^{-y^2}\,dy.
\]
当 $\lambda\to+\infty$ 时，积分上下限趋于 $-\infty,+\infty$。
利用高斯积分
\[
 \int_{-\infty}^{+\infty}e^{-y^2}\,dy=\sqrt\pi,
\]
便得到主要项
\[
 \boxed{
 F(\lambda)\sim
 f(x_0)e^{\lambda S(x_0)}
 \sqrt{\frac{2\pi}{-\lambda S''(x_0)}}.
 }
\]

直觉上，峰的高度约为 $e^{\lambda S(x_0)}$，主要贡献区域的宽度约为
$\lambda^{-1/2}$；$f$ 在这段越来越窄的区域里近似为 $f(x_0)$，
而峰的形状在换元后趋于高斯形状。

\subsection*{Watson 引理}
设 $\alpha,\beta>0$，$a\in(0,+\infty]$，$f:(0,a)\to\mathbb R$
连续且有界，$k$ 是非负整数，并且当 $x\to0^+$ 时，
\[
 f(x)=\sum_{p=0}^k a_px^p+o(x^k).
\]
则当 $\lambda\to+\infty$ 时，
\[
 \boxed{
 \int_0^a x^{\beta-1}f(x)e^{-\lambda x^\alpha}\,dx
 =\frac1\alpha\sum_{p=0}^k
 a_p\Gamma\!\left(\frac{\beta+p}{\alpha}\right)
 \lambda^{-(\beta+p)/\alpha}
 +o\!\left(\lambda^{-(\beta+k)/\alpha}\right).
 }
\]

\paragraph{证明.}
记 $F(\lambda)$ 为左侧积分。取充分小的固定正数 $b\in(0,a)$，
在 $(0,b)$ 上写
\[
 f(x)=\sum_{p=0}^k a_px^p+r(x),\qquad
 r(x)=x^kh(x),\qquad h(x)\longrightarrow0\quad(x\to0^+).
\]
取 $b$ 足够小，可保证 $h$ 在 $(0,b)$ 上有界。
用加法形式局部化，有
\[
 F(\lambda)=\sum_{p=0}^k a_p
 \int_0^b x^{\beta+p-1}e^{-\lambda x^\alpha}\,dx
 +R(\lambda)+O(e^{-\lambda b^\alpha}),
 \qquad
 R(\lambda)=\int_0^b x^{\beta-1}r(x)e^{-\lambda x^\alpha}\,dx.
\]

对每个 $p$，令 $u=\lambda x^\alpha$，则
\[
 x=(u/\lambda)^{1/\alpha},\qquad
 dx=\frac1\alpha\lambda^{-1/\alpha}u^{1/\alpha-1}\,du,
\]
故
\[
 \int_0^b x^{\beta+p-1}e^{-\lambda x^\alpha}\,dx
 =\frac1\alpha\lambda^{-(\beta+p)/\alpha}
 \int_0^{\lambda b^\alpha}u^{(\beta+p)/\alpha-1}e^{-u}\,du.
\]
对固定 $s>0$，由 Gamma 积分及尾部估计，
\[
 \int_0^\infty u^{s-1}e^{-u}\,du=\Gamma(s),\qquad
 \int_L^\infty u^{s-1}e^{-u}\,du
 \le e^{-L/2}\int_0^\infty u^{s-1}e^{-u/2}\,du
 =2^s\Gamma(s)e^{-L/2}.
\]
因此令 $s_p=(\beta+p)/\alpha$，得到
\[
 \int_0^b x^{\beta+p-1}e^{-\lambda x^\alpha}\,dx
 =\frac1\alpha\lambda^{-s_p}
 \left[\Gamma(s_p)+O(e^{-\lambda b^\alpha/2})\right].
\]

最后估计余项。令 $q=(\beta+k)/\alpha>0$，同样换元得
\[
 R(\lambda)=\frac1\alpha\lambda^{-q}
 \int_0^{\lambda b^\alpha}u^{q-1}e^{-u}
 h\!\left((u/\lambda)^{1/\alpha}\right)\,du.
\]
将积分中的函数在 $u\ge\lambda b^\alpha$ 处延拓为 $0$。
对每个固定 $u>0$，当 $\lambda\to+\infty$ 时，
$h((u/\lambda)^{1/\alpha})\to0$；且被积函数的绝对值不超过
\[
 C u^{q-1}e^{-u},\qquad
 \int_0^\infty C u^{q-1}e^{-u}\,du=C\Gamma(q)<\infty.
\]
由控制收敛定理，上面的积分趋于 $0$，因此
\[
 R(\lambda)=o(\lambda^{-q})
 =o\!\left(\lambda^{-(\beta+k)/\alpha}\right).
\]
所有指数小的误差也均为 $o(\lambda^{-q})$，合并即得结论。

\paragraph{RMK（端点权重与局部化的修补）.}
当 $0<\beta<1$ 时，$x^{\beta-1}$ 在 $0$ 处发散，
但 $f$ 有界且 $\beta>0$，所以
\[
 \int_0^b x^{\beta-1}|f(x)|\,dx
 \le\|f\|_\infty\frac{b^\beta}{\beta}<\infty.
\]
因此局部积分绝对收敛，不必要求振幅 $x^{\beta-1}f(x)$ 在端点连续。
局部化只需估计远离 $0$ 的尾部：固定 $\lambda_0>0$，对
$\lambda\ge\lambda_0$，有
\begin{align*}
 \left|\int_b^a x^{\beta-1}f(x)e^{-\lambda x^\alpha}\,dx\right|
 &\le e^{-(\lambda-\lambda_0)b^\alpha}
 \int_b^a x^{\beta-1}|f(x)|e^{-\lambda_0x^\alpha}\,dx\\
 &=O(e^{-\lambda b^\alpha}).
\end{align*}
这里的固定积分即使 $a=+\infty$ 也有限，因为它不超过
\[
 \frac{\|f\|_\infty}{\alpha}
 \lambda_0^{-\beta/\alpha}\Gamma(\beta/\alpha).
\]
这就直接给出了证明中所用的加法局部化；上述放缩使用的是绝对可积性。

\subsection*{用 Watson 引理推导 Laplace 方法}
设 $f,S\in C([a,b])$，$x_0\in(a,b)$ 是 $S$ 的唯一全局最大值点，
$S$ 在 $x_0$ 附近属于 $C^2$，且
\[
 S'(x_0)=0,\qquad A=-S''(x_0)>0.
\]
记 $M=S(x_0)$，考虑
\[
 F(\lambda)=\int_a^b f(x)e^{\lambda S(x)}\,dx.
\]

\paragraph{1. 局部化并精确化简相位.}
在 $x_0$ 的左侧，$S'(x)>0$；在右侧，$S'(x)<0$。
分别在两侧作换元
\[
 t=\sqrt{M-S(x)},\qquad S(x)=M-t^2.
\]
取充分小的固定 $\rho>0$，令两侧的反函数为
\[
 x=x_-(t)<x_0<x_+(t)=x,\qquad 0<t\le\rho,
 \qquad x_\pm(0)=x_0.
\]
于是 $U=(x_-(\rho),x_+(\rho))$ 是 $x_0$ 的固定邻域。
由加法局部化，存在 $\delta>0$，使
\[
 F(\lambda)=\int_U f(x)e^{\lambda S(x)}\,dx
 +O(e^{\lambda(M-\delta)}).
\]

\paragraph{2. 计算换元后的振幅.}
由 Taylor 展开，
\[
 M-S(x)=\frac A2(x-x_0)^2+o((x-x_0)^2),\qquad
 S'(x)=-A(x-x_0)+o(|x-x_0|).
\]
因此
\[
 \frac{x_\pm(t)-x_0}{t}\longrightarrow\pm\sqrt{\frac2A}.
\]
对 $S(x_\pm(t))=M-t^2$ 求导，得到
\[
 x_\pm'(t)=\frac{-2t}{S'(x_\pm(t))}
 \longrightarrow\pm\sqrt{\frac2A}.
\]
左支在 $x$ 从 $x_-(\rho)$ 增加到 $x_0$ 时，$t$ 从 $\rho$ 减少到 $0$。
合并两支，故
\[
 \int_U f(x)e^{\lambda S(x)}\,dx
 =e^{\lambda M}\int_0^\rho H(t)e^{-\lambda t^2}\,dt,
\]
其中
\[
 H(t)=f(x_+(t))x_+'(t)-f(x_-(t))x_-'(t).
\]
此函数可连续延拓到 $t=0$，并且
\[
 H(t)=H(0)+o(1),\qquad
 H(0)=2f(x_0)\sqrt{\frac2A}.
\]

\paragraph{3. 应用 Watson 引理.}
对上面的积分取 Watson 引理中的 $\alpha=2$、$\beta=1$、$k=0$，得
\begin{align*}
 \int_0^\rho H(t)e^{-\lambda t^2}\,dt
 &=\frac12H(0)\Gamma(1/2)\lambda^{-1/2}
   +o(\lambda^{-1/2})\\
 &=f(x_0)\sqrt{\frac{2\pi}{A}}\lambda^{-1/2}
   +o(\lambda^{-1/2}).
\end{align*}
再合并指数小的局部化误差，得到
\[
 \boxed{
 F(\lambda)=e^{\lambda S(x_0)}
 \left[f(x_0)\sqrt{\frac{2\pi}{-\lambda S''(x_0)}}
 +o(\lambda^{-1/2})\right].
 }
\]
若 $f(x_0)\ne0$，即可写成
\[
 F(\lambda)\sim f(x_0)e^{\lambda S(x_0)}
 \sqrt{\frac{2\pi}{-\lambda S''(x_0)}}.
\]
这里的换元使相位精确等于 $M-t^2$，而不是直接在指数中丢弃 Taylor 余项；
剩下的振幅渐近由 Watson 引理处理。

\end{document}
